% GENERATED FILE. Edit canonical lessons, then run npm run generate:books.
% canonical-source-hash: fnv1a64:f5dbbc5082fbf5d2
% canonical-lessons: GU-C349-vajbi, GU-C349-bhedbhav, GU-C349-saman, GU-C349-tatasth, GU-C349-vartvu

\chapter{Fair and Unfair}
\label{ch:gu-fair-and-unfair}
\hlchaptermodality{349}

\begin{chapteropening}
\textbf{By the end of this chapter:} \emph{I can say whether people are treated fairly.}

Complete the last lesson of chapter 349: I can say whether people are treated fairly.
\end{chapteropening}

\section[\texorpdfstring{\gu{વાજબી} (vājbī)}{વાજબી (vājbī)}]{\gu{વાજબી} (vājbī) — reasonable; fair}

\label{lesson:GU-C349-vajbi}



\begin{quote}
Before the new one: say the Gujarati for a ruling, then the Gujarati for to hit hard.
\end{quote}

\subsection*{You'll want to know: \gu{વાજબી}}
\textbf{\gu{વાજબી}} — \emph{vājbī} — \textquotedblleft{}reasonable; fair\textquotedblright{}.

\begin{quote}
\gu{દુકાનદારે} \gu{વાજબી} \gu{ભાવ} \gu{માગ્યો}.
\end{quote}

\begin{quote}
\emph{dukāndāre vājbī bhāv māgyo.}
\end{quote}

\begin{quote}
The shopkeeper asked a fair price.
\end{quote}

\begin{cousinweb}
\textbf{\gu{વાજબી}} (\emph{vājbī}) is usually traced to Arabic \emph{wājib}, \textquotedblleft{}due, proper\textquotedblright{}, through Persian.
\end{cousinweb}

\begin{grammarlens}[title={What you've built}]
One more word for giving reasons and opinions.
\end{grammarlens}

\subsection*{Your turn}
\begin{itemize}
\raggedright
  \item \emph{Say it:} \emph{vājbī}
  \item \emph{Say it:} \emph{vājbī}, once more
  \item \emph{Say it:} \emph{dukāndāre vājbī bhāv māgyo.} — The shopkeeper asked a fair price.
  \item \emph{Recall:} say the Gujarati for a ruling, then the Gujarati for to hit hard, then say \emph{vājbī} again
\end{itemize}

\begin{tcolorbox}[breakable,colback=teal!4,colframe=teal!35!black,title={Before you move on}]
What does \gu{વાજબી} mean? (\textquotedblleft{}Reasonable; fair\textquotedblright{}.) Say it once more.
\end{tcolorbox}
\section[\texorpdfstring{\gu{ભેદભાવ} (bhedbhāv)}{ભેદભાવ (bhedbhāv)}]{\gu{ભેદભાવ} (bhedbhāv) — discrimination}

\label{lesson:GU-C349-bhedbhav}



\begin{quote}
Before the new one: say the Gujarati for to hit hard, then the Gujarati for reasonable.
\end{quote}

\subsection*{You'll want to know: \gu{ભેદભાવ}}
\textbf{\gu{ભેદભાવ}} — \emph{bhedbhāv} — \textquotedblleft{}discrimination\textquotedblright{}.

\begin{quote}
\gu{શાળામાં} \gu{છોકરા} \gu{અને} \gu{છોકરી} \gu{વચ્ચે} \gu{ભેદભાવ} \gu{ન} \gu{હોવો} \gu{જોઈએ}.
\end{quote}

\begin{quote}
\emph{śāḷāmā̃ chokrā ane chokrī vacce bhedbhāv na hovo joīe.}
\end{quote}

\begin{quote}
There should be no discrimination between boys and girls at school.
\end{quote}

A masculine noun.

\begin{cousinweb}
\textbf{\gu{ભેદભાવ}} (\emph{bhedbhāv}) is built from Sanskrit \emph{bheda}, \textquotedblleft{}a difference\textquotedblright{}, and \emph{bhāva}, \textquotedblleft{}a state, a feeling\textquotedblright{} (Monier-Williams): treating people as different.
\end{cousinweb}

\begin{grammarlens}[title={What you've built}]
One more word for giving reasons and opinions.
\end{grammarlens}

\subsection*{Your turn}
\begin{itemize}
\raggedright
  \item \emph{Say it:} \emph{bhedbhāv}
  \item \emph{Say it:} \emph{bhedbhāv}, once more
  \item \emph{Say it:} \emph{śāḷāmā̃ chokrā ane chokrī vacce bhedbhāv na hovo joīe.} — There should be no discrimination between boys and girls at school.
  \item \emph{Recall:} say the Gujarati for to hit hard, then the Gujarati for reasonable, then say \emph{bhedbhāv} again
\end{itemize}

\begin{tcolorbox}[breakable,colback=teal!4,colframe=teal!35!black,title={Before you move on}]
What does \gu{ભેદભાવ} mean? (\textquotedblleft{}Discrimination\textquotedblright{}.) Say it once more.
\end{tcolorbox}
\section[\texorpdfstring{\gu{સમાન} (samān)}{સમાન (samān)}]{\gu{સમાન} (samān) — equal}

\label{lesson:GU-C349-saman}



\begin{quote}
Before the new one: say the Gujarati for reasonable, then the Gujarati for discrimination.
\end{quote}

\subsection*{You'll want to know: \gu{સમાન}}
\textbf{\gu{સમાન}} — \emph{samān} — \textquotedblleft{}equal\textquotedblright{}.

\begin{quote}
\gu{કાયદા} \gu{સામે} \gu{બધા} \gu{માણસો} \gu{સમાન} \gu{છે}.
\end{quote}

\begin{quote}
\emph{kāydā sāme badhā māṇso samān che.}
\end{quote}

\begin{quote}
All people are equal before the law.
\end{quote}

\begin{cousinweb}
\textbf{\gu{સમાન}} (\emph{samān}) is a tatsama, borrowed from Sanskrit \emph{samāna}, \textquotedblleft{}same, alike, equal\textquotedblright{} (Monier-Williams).
\end{cousinweb}

\begin{grammarlens}[title={What you've built}]
One more word for giving reasons and opinions.
\end{grammarlens}

\subsection*{Your turn}
\begin{itemize}
\raggedright
  \item \emph{Say it:} \emph{samān}
  \item \emph{Say it:} \emph{samān}, once more
  \item \emph{Say it:} \emph{kāydā sāme badhā māṇso samān che.} — All people are equal before the law.
  \item \emph{Recall:} say the Gujarati for reasonable, then the Gujarati for discrimination, then say \emph{samān} again
\end{itemize}

\begin{tcolorbox}[breakable,colback=teal!4,colframe=teal!35!black,title={Before you move on}]
What does \gu{સમાન} mean? (\textquotedblleft{}Equal\textquotedblright{}.) Say it once more.
\end{tcolorbox}
\section[\texorpdfstring{\gu{તટસ્થ} (taṭasth)}{તટસ્થ (taṭasth)}]{\gu{તટસ્થ} (taṭasth) — neutral}

\label{lesson:GU-C349-tatasth}



\begin{quote}
Before the new one: say the Gujarati for discrimination, then the Gujarati for equal.
\end{quote}

\subsection*{You'll want to know: \gu{તટસ્થ}}
\textbf{\gu{તટસ્થ}} — \emph{taṭasth} — \textquotedblleft{}neutral\textquotedblright{}.

\begin{quote}
\gu{આ} \gu{વિવાદમાં} \gu{હું} \gu{તટસ્થ} \gu{રહીશ}.
\end{quote}

\begin{quote}
\emph{ā vivādmā̃ hũ taṭasth rahīś.}
\end{quote}

\begin{quote}
I will stay neutral in this dispute.
\end{quote}

\begin{cousinweb}
\textbf{\gu{તટસ્થ}} (\emph{taṭasth}) is a tatsama, borrowed from Sanskrit \emph{taṭastha}, \textquotedblleft{}standing on the bank\textquotedblright{} (Monier-Williams): one who watches the river from the bank. Gujarati uses it for taking neither side.
\end{cousinweb}

\begin{grammarlens}[title={What you've built}]
One more word for giving reasons and opinions.
\end{grammarlens}

\subsection*{Your turn}
\begin{itemize}
\raggedright
  \item \emph{Say it:} \emph{taṭasth}
  \item \emph{Say it:} \emph{taṭasth}, once more
  \item \emph{Say it:} \emph{ā vivādmā̃ hũ taṭasth rahīś.} — I will stay neutral in this dispute.
  \item \emph{Recall:} say the Gujarati for discrimination, then the Gujarati for equal, then say \emph{taṭasth} again
\end{itemize}

\begin{tcolorbox}[breakable,colback=teal!4,colframe=teal!35!black,title={Before you move on}]
What does \gu{તટસ્થ} mean? (\textquotedblleft{}Neutral\textquotedblright{}.) Say it once more.
\end{tcolorbox}
\section[\texorpdfstring{\gu{વર્તવું} (vartvũ)}{વર્તવું (vartvũ)}]{\gu{વર્તવું} (vartvũ) — to behave}

\label{lesson:GU-C349-vartvu}



\begin{quote}
Before the new one: say the Gujarati for equal, then the Gujarati for neutral.
\end{quote}

\subsection*{You'll want to know: \gu{વર્તવું}}
\textbf{\gu{વર્તવું}} — \emph{vartvũ} — \textquotedblleft{}to behave\textquotedblright{}.

\begin{quote}
\gu{તમારે} \gu{બધા} \gu{સાથે} \gu{સમાન} \gu{રીતે} \gu{વર્તવું} \gu{જોઈએ}.
\end{quote}

\begin{quote}
\emph{tamāre badhā sāthe samān rīte vartvũ joīe.}
\end{quote}

\begin{quote}
You should treat everyone equally.
\end{quote}

Past (he …): \gu{તે} \gu{મહેમાનો} \gu{સાથે} \gu{સારી} \gu{રીતે} \gu{વર્ત્યો}, \emph{te mahemāno sāthe sārī rīte vartyo}.

\begin{cousinweb}
\textbf{\gu{વર્તવું}} (\emph{vartvũ}) goes back to Sanskrit \emph{vártatē}, \textquotedblleft{}turns, moves; takes place\textquotedblright{} (CDIAL 11352). CDIAL lists no Gujarati form; \textquotedblleft{}to behave\textquotedblright{}, to carry oneself, is the sense Gujarati gives it.
\end{cousinweb}

\begin{grammarlens}[title={What you've built}]
One more word for giving reasons and opinions.
\end{grammarlens}

\subsection*{Your turn}
\begin{itemize}
\raggedright
  \item \emph{Say it:} \emph{vartvũ}
  \item \emph{Say it:} \emph{vartvũ}, once more
  \item \emph{Say it:} \emph{tamāre badhā sāthe samān rīte vartvũ joīe.} — You should treat everyone equally.
  \item \emph{Recall:} say the Gujarati for equal, then the Gujarati for neutral, then say \emph{vartvũ} again
\end{itemize}

\begin{tcolorbox}[breakable,colback=teal!4,colframe=teal!35!black,title={Before you move on}]
What does \gu{વર્તવું} mean? (\textquotedblleft{}To behave\textquotedblright{}.) Say it once more.
\end{tcolorbox}
